%% Copyright 2026 Petr Socha
%%
%% This work may be distributed and/or modified under the
%% conditions of the LaTeX Project Public License, either version 1.3c
%% of this license or (at your option) any later version.
%% The latest version of this license is in
%%   https://www.latex-project.org/lppl.txt
%% and version 1.3c is part of all distributions of LaTeX
%% version 2008 or later.
%%
%% This work has the LPPL maintenance status `maintained'.
%% The Current Maintainer of this work is Petr Socha.
%% Maintenance and bug reports: https://github.com/petrsocha/pslib-latex/issues
%%
%% The components of this work are listed in MANIFEST.txt.
%% The school logo is a separately licensed third-party asset.

% ============================================================================
% NÁVRH -- OBRÁZEK, TABULKA, POPISKY A ODKAZY
% ============================================================================
\chapter{Návrh}
\label{chap:navrh}

% Následující odstavce jsou pouze ilustrační výplňový text.
% Při psaní vlastní práce je nahraďte popisem navrženého řešení.
V~této kapitole je uvedena jednoduchá ukázka technického návrhu.
Text slouží především k~předvedení sazby obrázků, tabulek a~odkazů.
Ve skutečné práci by měl čtenář zjistit, jaké řešení bylo zvoleno,
proč je vhodné a~jak spolu jeho jednotlivé části souvisejí.

\section{Ukázkový obrázek}

Nejprve je vhodné popsat hlavní části navrhovaného systému a~jejich
vzájemné vazby. Přehledné schéma pomůže čtenáři pochopit princip řešení
dříve, než začne číst podrobnosti o~jednotlivých součástkách.
Přenos dat od senzoru znázorňuje \autoref{fig:tok-dat}:
mikrokontrolér předává naměřené hodnoty do počítače.

% Obrázky ukládejte do složky obrazky/ (PNG, JPG nebo PDF).
% Prostředí figure je plovoucí: [htbp] preferuje začátek/konec stránky,
% případně samostatnou stránku s obrázkem. Nenuťte obrázek na pevné místo.
% Popisek je POD obrázkem. Značku \label vložte ZA \caption.
% Odkaz v textu vložte pomocí \autoref{fig:tok-dat}.
\begin{figure}[htbp]
  \centering
  \includegraphics[width=0.78\textwidth]{obrazky/ukazkove-schema.pdf}
  \caption{Ukázkové schéma toku dat (vlastní zpracování)}
  \label{fig:tok-dat}
\end{figure}

Senzor získává měřenou veličinu, mikrokontrolér zpracuje přijaté údaje
a~počítač je zobrazí. Ve skutečné práci je vhodné připojit i~zdůvodnění
výběru jednotlivých součástek. Jejich parametry lze stručně porovnat
v~tabulce, která bývá přehlednější než souvislý výčet hodnot.

\section{Ukázková tabulka}

% Hodnoty jsou převzaty z oficiální dokumentace Arduino UNO R3.
% Tabulka má popisek NAD tabulkou a zarovnání pomocí prostředí tabular.
Dokumentace výrobce uvádí u~desky Arduino UNO R3 například procesor
ATmega328P a~takt 16\,MHz~\cite{arduinoUNO}.
Vybrané parametry shrnuje \autoref{tab:arduino-parametry}.

\begin{table}[htbp]
  \centering
  \caption{Vybrané parametry desky Arduino UNO R3 podle výrobce~\cite{arduinoUNO}}
  \label{tab:arduino-parametry}
  \begin{tabular}{ll}
    \hline
    Parametr & Hodnota \\
    \hline
    Mikrokontrolér & ATmega328P \\
    Taktovací frekvence & 16\,MHz \\
    Digitální vstupy/výstupy & 14 \\
    Analogové vstupy & 6 \\
    \hline
  \end{tabular}
\end{table}

Podrobnější parametry bývají uvedené v~datasheetu.
Jedním z~příkladů je technický list výrobce
STMicroelectronics~\cite{stm32datasheet}. Při porovnávání je nutné
ověřit jednotky a~konkrétní variantu součástky.

% U převzatých obrázků a údajů uveďte zdroj. U vlastních napište
% „vlastní zpracování“. Tabulky nepořizujte jako screenshoty.
